\section{配套合作服务与认证评审}

在冷弯薄壁型钢有限元分析模块开发、交付及后续 BP 软件全生命周期迭代过程中，华侨大学同步为大禾众邦提供以下配套合作服务，并牵头组织软件认证评审工作。

\subsection{四项专属配套服务}

\subsubsection{BP软件现有计算功能全面验证服务}
华侨大学组织冷弯钢专业研究团队，针对企业已完成开发上线的 BP v4.0 软件幕墙计算、整屋计算等两大功能开展逐项核验，出具完整 BP 软件功能验证记录与验证报告。

\subsubsection{BP软件高校专业课程培训服务}
依托华侨大学土木工程学院教学资源，协助大禾众邦企业面向轻钢设计行业从业者提供软件培训课程，配套制作教学课件、培训指导手册，开展软件实操教学，搭建企业人才培养渠道。

\subsubsection{BP软件 \(^+\) 有限元模块全方位联合测试服务}
项目全周期内，华侨大学测试团队全程配合企业研发人员开展联合测试：覆盖模型互通、静力/稳定/模态计算、多规范切换、大型工程模型压力测试、异常边界容错测试等场景；完整记录软件缺陷、问题跟踪整改，形成闭环测试台账。

\subsubsection{软件版本迭代认证评审组织服务}
本次有限元模块开发交付，以及后续 BP 软件重大版本升级迭代时，统一由华侨大学牵头组织全套软件认证评审工作；评审、第三方测评过程产生的全部费用均由大禾众邦承担。

\subsection{软件认证评审实施流程}
针对 BP 软件及新增有限元分析模块，设置专家评审 + CNAS 第三方独立测评两级验证流程，为软件工程应用、市场推广提供权威技术佐证。

\subsubsection{专家评审}
\begin{enumerate}[label=\arabic*., leftmargin=2em]
    \item 组织方：华侨大学；
    \item 评审专家组构成：高校钢结构/有限元领域教授、甲级设计院资深结构工程师；
    \item 评审内容：针对有限元底层单元算法、静力/稳定/模态计算逻辑、GB50018/AISI S100/EC3 多国规范条文代码实现、构件承载力与屈曲输出结果合理性开展算例专项核查；
    \item 交付成果：正式专家评审意见书。
\end{enumerate}

\subsubsection{CNAS 资质第三方独立测评}
\begin{enumerate}[label=\arabic*., leftmargin=2em]
    \item 组织对接方：华侨大学；
    \item 实施主体：具备CNAS检测资质的第三方工程软件测评机构；
    \item 测评内容：提供全套标准基准解析算例、构件试验模型、轻钢整体工程模型，由第三方独立复算，核验软件计算精度、稳定性、规范符合性；
    \item 交付成果：具备公信力的第三方正式测评报告。
\end{enumerate}

补充说明：国内建筑结构设计软件无强制准入认证，但行业惯例要求软件新增核心计算模块、重大版本更新时完成上述两级评审测评，相关报告可用于设计院工程采信、产品市场宣传、项目投标技术支撑。

\subsection{权责约定}
\begin{enumerate}[label=\arabic*., leftmargin=2em]
    \item 华侨大学职责：负责四项配套服务落地、专家团队邀约、第三方测评机构对接、组织评审全流程工作；
    \item 大禾众邦职责：提供BP软件安装包、工程模型、企业技术资料；承担专家劳务、第三方检测、会务、试验等全部评审测评相关费用；配合完成软件测试、教学素材供给。
\end{enumerate}